% Pursuing Stacks (cotes 134-2 to 134-8): archive cote 134-7, site page 54 (PDF page 55), typescript page 496.
% typescript. Draft compared: ps6.tex:2663-2741.
% AI-transcribed by Claude Opus 5.5 (claude-opus-5-5, adaptive thinking, effort high), Message Batches API, batch ps4 (msgbatch_012Aoc1NuskATPsyBeKSYK7J), 2026-11-12; not yet checked by eye.
\scanpage{134-7/54}
\mspage{496}

basic structures as $K(\underline{Z},n)$ with $n\geqslant 2$. At any rate, the \hand{``}way-out\hand{''} argument
I have finally been thinking of, in order to try and check (6) is a quasi%
isomorphism, rests on vanishing assumptions which (as I was informed by
Illusie the next day) are wholly unrealistic. This finally clears up,
it would seem, a tenacious misconception which has been sticking to my
first heuristic ponderings about the homology and cohomology formalism
for schematic homotopy types: \uline{one should be very careful not to substitu\xout{t}}%
\uline{te the ``pointed'' linearization} $L_{pt}(W(M))\simeq\Gamma_k(M)$ \uline{for the nonpointed}
\uline{one} $L(W(M))\simeq\Gamma_k{}^{\wedge}(M)$, \uline{in computing homology and cohomology invariants}
\uline{of schematic homotopy types}. To say it differently, in order to be
able to \xout{define} compute (or just define) the ``schematic'' homology and
cohomology invariants, we do need as a model a fullfledged semisimplic\unsure{i}\hand{al}
unipotent bundle, not just the corresponding formal one, giving rise
to invariants of \deleted{their}\inserted{it's} own, namely ``formal'' homology and cohomology,
which are definitely distinct from the former.

\deleted{Itxisxall thexmorexremarkablexthat}
118. It is all the more remarkable, in view of the preceding findings,
that the homotopy invariants
\[ \hand{\pi}_i(X_*) \ , \quad i \hand{\geqslant} 0 , \]
of a \inserted{pointed} semisimplicial \inserted{unipotent} bundle $X_*$ (still assuming the components $X_n$ to be
smooth, i.e. \inserted{to} correspond\deleted{ing} to projective modules) turn out to be invari-%
ants of the corresponding ``formal'' object, and, more startling still,
of the corresponding ``infinitesimal'' object of order 1. M\over{\unsure{g}}{o}re specifical-%
ly, consider the ``Lie functor'' or ``tangent space at the origin''
\[ (7) \qquad \mathrm{Lie}: U(k)^{\bullet} \longrightarrow Ab_k \ , \qquad X = W(M) \longmapsto M \ , \]
which we'll need only for the time being for smooth $X$, when the geometri
meaning of it is clear. This functor transforms semisimplicial \inserted{pointed} bundles
into semisimplicial $k$-modules $\mathrm{Lie}(X_*)$ , thus we should get, besides
abelianization, another remarkable functor\hand{,} from $\hand{\mathcal{M}_1}(k)$ \inserted{say} to $D.(Ab_k)$ :
\[ (8) \qquad \mathrm{Lie}: \hand{\mathcal{M}_1}(k) \longrightarrow D.(Ab_k) \ \inserted{via}, \ X_* \longmapsto \mathrm{Lie}(X_*) \ , \]
granting that $X_* \mapsto \mathrm{Lie}(X_*)$ transforms quasiisomorphisms into quasiiso-%
morphisms. Now that this must be so follows from the really startling
formula
\[ (9) \qquad \hand{\pi}_i(X_*) \hand{\simeq} \hand{\pi}_i(\mathrm{Lie}(X_*)) \quad (\hand{\simeq} H_i \text{ of the associated chain} \]
\[ \text{complex in } Ab_k \ ) \ , \]
where the left hand side, I recall, is defined as
\[ \hand{(10)} \qquad \hand{\pi}_i(X_*) = \hand{\pi}_i(X_*(k)) \ . \]

\margin{54}
